\section{Related Work}
\label{sec:related_work}

\paragraph{Static Parameter-Space Model Merging.}
Fusing multiple deep models fine-tuned from a shared pre-trained initialization has emerged as a key area of study to bypass expensive multi-task training. \citet{wortsman2022model} introduced Model Soups, showing that averaging the weights of models fine-tuned under different hyperparameter configurations improves generalization. To merge models trained on distinct downstream tasks, \citet{ilharco2022editing} proposed Task Arithmetic, which adds or subtracts task vectors (differences between fine-tuned and pre-trained weights) to steer a model's behavior. To mitigate the issue of sign conflicts and redundant parameter values, TIES-Merging \cite{yadav2023ties} trims negligible weights, resolves sign disagreements, and averages the remaining parameters. Other static merging methods employ Fisher information matrices \cite{matena2022merging} or covariance structures \cite{adamerging} to align representations. While these static fusions require zero inference-time overhead, they struggle on highly dissimilar, high-conflict target tasks due to structural representation conflicts, where a single static set of weights cannot accommodate divergent feature spaces.

\paragraph{Dynamic Test-Time Model Merging.}
To overcome the representation limits of static merging, several recent methods propose dynamic, test-time parameter ensembling. Quantum Wavefunction Superposition (QWS-Merge) \cite{qwsmerge} uses a wave-like interference projection to fuse weights. PolyMerge \cite{polymerge} fits polynomial trajectories across the coefficient space at test-time, while ZipMerge \cite{zipmerge} jointly optimizes pruning masks and merging parameters. Other dynamic frameworks include SuiteMerge \cite{suitemerge} and OFS-Tune \cite{ofstune}. 
While these methods achieve high multi-task performance by dynamically adapting weights to the input data, they ignore the severe \textbf{real-world deployment bottleneck}: performing high-dimensional weight-matrix reconstruction online. This on-the-fly reconstruction introduces huge computational and memory-transfer latency during inference. In contrast, our Hybrid-Router takes a highly practical engineering approach, utilizing layer-wise partitioning to compress the active parameter ensembling space to a small fraction of the network, drastically cutting latency.

\paragraph{Parameter-Efficient Fine-Tuning (PEFT) and Adapter Serving.}
PEFT methods, such as LoRA \cite{hu2021lora}, Bottleneck Adapters \cite{houlsby2019parameter}, and Prefix Tuning \cite{li2021prefix}, freeze the pre-trained backbone and inject a small set of trainable parameters. Multi-task serving via task-specific LoRA adapters over a shared base model is a popular industry standard, supported by specialized dynamic serving runtimes like Punica \cite{punica}, S-LoRA \cite{slora}, dLoRA \cite{dlora}, or Decoupled PEFT \cite{decoupled_peft}. However, dynamic adapter serving suffers from three major limitations: (1) \emph{Inference Engine and Hardware Lock-in:} These runtimes rely on highly complex, hardware-dependent custom CUDA/Triton kernels to execute batched GEMM operations across heterogeneous adapters. They cannot be easily compiled or executed on standard commodity edge hardware, microcontrollers, or web browsers. (2) \emph{Active Parameter VRAM Footprint:} Serving $K$ active adapters requires keeping all $K$ sets of adapter weights active in memory simultaneously. (3) \emph{Compatibility:} In contrast, Hybrid-Router operates strictly via standard parameter-blending operations. Once the weights are reconstructed, the model can be compiled and executed on *any* standard lightweight inference engine (TensorRT, ONNX Runtime, TFLite, CoreML) without custom kernel dependencies. Additionally, our layer partition framework reduces active task-vector parameters stored in VRAM by exactly $1 - k/L$ (e.g., a \textbf{71.4\%} savings at $k=4$), enabling dynamic multi-task intelligence on highly resource-constrained edge hardware.

\paragraph{Mixture of Experts (MoE) and Routing.}
Standard Mixture-of-Experts architectures \cite{jacobs1991adaptive, shazeer2017outrageously, fedus2021switch} employ routing layers that dynamically route individual tokens or sequence segments to specialized feedforward sub-networks. However, MoEs must be trained from scratch using complex auxiliary load-balancing losses, and their total parameter size is extremely large, making them difficult to deploy on consumer hardware. Our Hybrid-Router, in contrast, is a post-hoc, calibration-based ensembling method. It operates on independently fine-tuned specialists and requires only a tiny calibration set (64 samples) to optimize a lightweight, 772-parameter routing head, making it extremely cost-effective and easy to integrate into existing production systems.
